% !PRINTED 16
\weberpage{16}

%% [running head, not body text -- commented out 2026-08-07] Page furniture (running
%% head, folio, signature) is excluded from this edition by standing policy.
%% 16 \hfill \gesperrt{Erster Abschnitt.} \hfill \gesperrt{\S.\,5.}

%<<CONT
the isomorphic groups (\S.\,1), only with the difference that to each element $A$ there corresponds not merely one element $a$, but several. This fact gives occasion to extend the concept of isomorphism, as is done by the following definition:

\gesperrt{A group $P$ with the elements $a, b, c \dots$ is called (multiply) isomorphic with a group $Q$ with the elements $A, B, C \dots$ if the two groups can be so related to one another that to each of the elements $A, B, C \dots$ there correspond one or more of the elements $a, b, c \dots$, and indeed in such a way that each of the elements of $Q$ corresponds to one and only one of the elements of $P$, and that, when $a$ and $A$, $b$ and $B$ correspond to one another, $ab$ corresponds to the element $AB$.}
% [sic: Weber prints ``elements of Q''; the uniqueness condition requires P.]

Thus it can first be proved that to each of the elements $A, B, C \dots$ there corresponds an equal number of elements $a, b, c \dots$, and that therefore \gesperrt{the degree of $Q$ is a divisor of the degree of $P$}. For let $A$ be the unit element in $Q$, to which the elements $a, a_1 \dots a_{\nu-1}$ in $P$ may correspond. These latter elements must then form a group of degree $\nu$, which we shall denote by $N$, because by the definition of isomorphism the element $a a_1$ must correspond to the element $AA = A$. If then $b$ is an element of $P$ corresponding to the element $B$, then again all the elements $Nb$ must correspond to the same element $B$ of $Q$. For every $ab$ must correspond to the element $AB = B$. If $b_1, b$ are two elements corresponding to the element $B$, then $b_1 b^{-1} = a$ corresponds to the element $A$, and is therefore contained in $N$; therefore $b_1 = ab$ is contained in $Nb$. By $Nb$, therefore, all the elements that correspond to the element $B$ are exhausted, and to each element of $Q$ there correspond $\nu$ elements of $P$. The isomorphism is called \gesperrt{$\nu$-fold}. But every element $b^{-1}ab$ likewise corresponds to the unit element $A$, and therefore $b^{-1}Nb = N$, i.\,e. $N$ is a normal divisor of $P$, and $Q$ is simply isomorphic with $P/N$.

We see, then, that there is no other multiple isomorphism than that between the complementary group of a normal divisor and the whole group. But one could extend the concept of isomorphism still further, by calling two groups which are $\mu$- and $\nu$-fold
%CONT>>
